% !TeX program = tectonic
\documentclass[11pt,a4paper]{article}
\usepackage{fontspec}
\usepackage[english]{babel}
\babelfont{rm}[Language=Default]{Liberation Serif}
\babelfont[Language=Default]{sf}{Carlito}
\babelfont[Language=Default]{tt}{DejaVu Sans Mono}
\usepackage{amsmath,amssymb,amsthm}
\usepackage{geometry}
\geometry{a4paper, top=2.4cm, bottom=2.4cm, left=2.4cm, right=2.4cm}
\usepackage{booktabs,tabularx,enumitem,xcolor,tcolorbox}
\tcbuselibrary{breakable,skins}
\definecolor{linkblue}{HTML}{1F4E79}
\definecolor{statgreen}{HTML}{2F6B3A}
\definecolor{goldacc}{HTML}{8B7E5A}
\usepackage[colorlinks=true,linkcolor=linkblue,urlcolor=linkblue,unicode=true]{hyperref}
\theoremstyle{plain}
\newtheorem{theorem}{Theorem}
\newtheorem{lemma}[theorem]{Lemma}
\newtheorem{proposition}[theorem]{Proposition}
\newtheorem{corollary}[theorem]{Corollary}
\theoremstyle{definition}
\newtheorem{definition}[theorem]{Definition}
\theoremstyle{remark}
\newtheorem{remark}[theorem]{Remark}
\newtcolorbox{statusbox}[1][]{colback=statgreen!5,colframe=statgreen!75,
  title={\textbf{Summary of results:} #1},fonttitle=\small,boxrule=0.7pt,arc=2pt,breakable}
\newtcolorbox{databox}[1][]{colback=goldacc!6,colframe=goldacc!80,
  title={\textbf{Protocol data:} #1},fonttitle=\small,boxrule=0.7pt,arc=2pt,breakable}
\newtcolorbox{verifybox}[1][]{colback=linkblue!5,colframe=linkblue!80,
  title={\textbf{Verification:} #1},fonttitle=\small,boxrule=0.7pt,arc=2pt,breakable}
\widowpenalty=10000
\clubpenalty=10000
\setlength{\parskip}{0.35em}
\newcommand{\CC}{\mathbb{C}}
\newcommand{\RR}{\mathbb{R}}
\newcommand{\ZZ}{\mathbb{Z}}
\newcommand{\QQ}{\mathbb{Q}}
\newcommand{\Ga}{\Gamma}
\newcommand{\im}{\operatorname{Im}}
\newcommand{\re}{\operatorname{Re}}
\newcommand{\disc}{\operatorname{disc}}
\begin{document}
\begin{center}
{\LARGE\bfseries The Quartic Tower: Exact Radicals}\\[0.2em]
{\large and the Minimal Quartics of the Rungs $N=15/30$}\\[0.6em]
{\normalsize Isaev Iskhak Khamzatovich}\\[0.2em]
{\small The ``Dynamic Principle'' program --- the \texttt{hodge-laboratory} repository}\\[0.15em]
{\small Standalone theorem monograph · Monograph 19/20}
\end{center}
\vspace{0.4em}
{\color{linkblue}\hrule height 0.8pt}
\vspace{0.8em}

\section{Setting}

The levels $N=15$ and $N=30$ are the constructible rungs of the
ladder: $15=3\cdot5$ is a product of distinct Fermat primes, and
the regular pentadecagon is constructible by straightedge and
compass (Gauss--Wantzel). Accordingly $2\cos(2\pi/15)$ has degree
$4$ over $\QQ$ ($\varphi(15)/2=4$), and $b_{\mathrm{Ch}}(15)$,
$b_{\mathrm{Ch}}(30)$ are expressible in square radicals --- by a
two-storey tower $\QQ\subset\QQ(\sqrt5)\subset\QQ(\sqrt5,\beta)$.
This monograph fixes the complete exact layer of v1.2: the closed
radicals, the minimal quartics, the Galois structure of the
tower, and the machine certificate --- the integer identity check
in $\ZZ[\sqrt5][\beta]$.

\begin{lemma}[doubling permutes the classes]\label{th:dbl}
The set $C=\{1,2,4,7\}$ contains one representative of each pair
of classes $\{\pm k\}$ of invertible residues modulo $15$.
Multiplication by $2$ permutes $C$ (up to changing the
representative): $1\mapsto2\mapsto4\mapsto8\equiv-7\mapsto7\mapsto14\equiv-1\mapsto1$.
\end{lemma}

\begin{proof}
Direct check: $2\cdot1=2$, $2\cdot2=4$, $2\cdot4=8\equiv8$, and
$8\equiv-7\pmod{15}$; further $2\cdot7=14\equiv-1$. The classes
$\{\pm1\},\{\pm2\},\{\pm4\},\{\pm7\}$ are cycled into one another
--- this is precisely the cyclic group
$\mathrm{Gal}(\QQ(\zeta_{15})^+/\QQ)
\cong(\ZZ/15)^\times/\{\pm1\}\cong C_4$ generated by doubling.
\end{proof}

\begin{lemma}[Ramanujan sums]\label{th:rams}
For $x_k=\zeta^k+\zeta^{-k}$, $\zeta=e^{2\pi i/15}$, $k\in C$:
$s_1=\sum x_k=1$ and $s_2=\sum_{i<j}x_ix_j=-4$.
\end{lemma}

\begin{proof}
$s_1=\sum_{k\in C}(\zeta^k+\zeta^{-k})=\sum_{j\in U}\zeta^j$,
where $U=(\ZZ/15)^\times$ runs over all eight units; this is the
Ramanujan sum $c_{15}(1)=\mu(15)=\mu(3)\mu(5)=1$. Next,
$x_k^2=2+2\cos(4\pi k/15)=2+x_{2k}$, so by
Lemma~\ref{th:dbl} $\sum_{k\in C}x_k^2=8+s_1=9$, and
$s_2=\frac{s_1^2-\sum x_k^2}{2}=\frac{1-9}{2}=-4$.
\end{proof}

\begin{theorem}[the tower and the radicals]\label{th:tower}
Put $\beta=\sqrt{30-6\sqrt5}$ and $\beta'=\sqrt{30+6\sqrt5}$
(positive roots). Then
\[
  2\cos\frac{2\pi}{15}=\frac{1+\sqrt5+\beta}{4},
  \qquad
  2\cos\frac{2\pi}{30}=\frac{\sqrt5-1+\beta'}{4},
\]
and $\beta\beta'=12\sqrt5$ is an exact identity, so that
$\beta'=12\sqrt5/\beta$ and both radicals live in \emph{one}
quartic field $K=\QQ(\sqrt5,\beta)=\QQ(\sqrt5,\beta')
=\QQ(\zeta_{15})^+=\QQ(\zeta_{30})^+$.
The field $K$ is totally real and cyclic:
$\mathrm{Gal}(K/\QQ)\cong C_4$, and $\QQ(\sqrt5)$ is its unique
quadratic subfield.
\end{theorem}

\begin{proof}
\emph{The forms.} For $x=(1+\sqrt5+\beta)/4$ we have
$(4x-1-\sqrt5)^2=\beta^2=30-6\sqrt5$; expanding,
$16x^2-8(1+\sqrt5)x+(6+2\sqrt5)=30-6\sqrt5$, i.e.
$x^2=\varphi\,x+\gamma$ with $\varphi=\frac{1+\sqrt5}{2}$ and
$\gamma=\frac{3-\sqrt5}{2}$: $x$ is quadratic over $\QQ(\sqrt5)$.
Moreover $x\notin\QQ(\sqrt5)$: otherwise the quadratic
$z^2-\varphi z-\gamma$ would have a root in $\QQ(\sqrt5)$, and its
discriminant $\Delta_z=\varphi^2+4\gamma=\frac{3(5-\sqrt5)}{2}$
would have norm $\mathrm{N}(\Delta_z)=\frac94\cdot20=45$, a
rational square ($\mathrm{N}(w^2)=\mathrm{N}(w)^2\in\QQ^2$) ---
but $45$ is not a square. Hence
$[\QQ(\sqrt5,x):\QQ(\sqrt5)]=2$ and $[\QQ(\sqrt5,x):\QQ]=4$.
Since $x=2\cos(2\pi/15)$ lies in $\QQ(\zeta_{15})^+$ and
$\varphi(15)/2=4$, the fields coincide: $K=\QQ(\zeta_{15})^+$.
For the second form: $\zeta_{30}=-\zeta_{15}^8$ (an order check),
so $\QQ(\zeta_{30})=\QQ(\zeta_{15})$, and the analogous expansion
$(4x'-\sqrt5+1)^2=\beta'^2$ gives
$x'=(\sqrt5-1+\beta')/4=2\cos(\pi/15)=2\cos(2\pi/30)$.

\emph{The tower pairing.}
$\beta^2\beta'^2=(30-6\sqrt5)(30+6\sqrt5)=900-180=720=(12\sqrt5)^2$,
and $\beta,\beta'>0$; therefore $\beta\beta'=12\sqrt5$ exactly,
whence $\beta'=12\sqrt5/\beta\in K$.

\emph{Cyclicity.} $\mathrm{Gal}(\QQ(\zeta_{15})^+/\QQ)
\cong(\ZZ/15)^\times/\{\pm1\}$; the element $2$ has order $4$
modulo $15$, so the group is cyclic of order $4$.
The automorphism $\sigma\colon\zeta\mapsto\zeta^2$ sends
$\sqrt5=\zeta_5+\zeta_5^4-\zeta_5^2-\zeta_5^3$ (where
$\zeta_5=\zeta^3$) to $-\sqrt5$; hence
$\sigma(\beta)^2=\sigma(\beta^2)=30+6\sqrt5=\beta'^2$ and, by the
pairing, $\sigma(\beta)=\pm\beta'$ --- the sign is pinned:
$\sigma$ swaps $\beta\leftrightarrow\beta'$ up to sign, while
$\langle\sigma^2\rangle$ fixes $\QQ(\sqrt5)$ --- the unique
quadratic subfield (the group $C_4$ has a unique subgroup of
order $2$). The discriminant of the field: by the
conductor--discriminant formula
$\disc K=f_4\cdot f'_4\cdot f_2=15\cdot15\cdot5=1125=3^2\cdot5^3$
--- the very number that equals the volume
$\mathrm{vol}_h=1125$ of monograph 16/20.
\end{proof}

\begin{theorem}[the minimal quartics]\label{th:quart}
The minimal polynomials are
\[
  P_{15}(x)=x^4-x^3-4x^2+4x+1 \ \text{for}\ 2\cos\frac{2\pi}{15},
  \qquad
  P_{30}(x)=x^4+x^3-4x^2-4x+1 \ \text{for}\ 2\cos\frac{2\pi}{30},
\]
and $P_{30}(x)=P_{15}(-x)$ --- the level-$30$ quartic is the
reflection of the level-$15$ quartic. The coefficients of
$P_{15}$ assemble from Ramanujan sums and the tower norm: the
constant term $s_4=\prod x_k=1$ and the third coefficient
$s_3=-4$ (Newton).
\end{theorem}

\begin{proof}
The quartic with roots $x_1,\dots,x_4$ (the classes $C$) reads
$x^4-s_1x^3+s_2x^2-s_3x+s_4$. By Lemma~\ref{th:rams} already
$s_1=1$, $s_2=-4$. \emph{The tower norm}: by
Theorem~\ref{th:tower} $K\supset\QQ(\sqrt5)$ and the norm splits:
$x_1x_3=(\zeta+\zeta^{-1})(\zeta^4+\zeta^{-4})
=(\zeta^5+\zeta^{-5})+(\zeta^3+\zeta^{-3})
=-1+2\cos\frac{2\pi}{5}=-1+\frac{\sqrt5-1}{2}=\frac{\sqrt5-3}{2}$;
then
$s_4=\mathrm{N}_{K/\QQ}(x_1)
=\mathrm{N}_{\QQ(\sqrt5)/\QQ}\!\left(\frac{\sqrt5-3}{2}\right)
=\frac{\sqrt5-3}{2}\cdot\frac{-\sqrt5-3}{2}=1$.
\emph{Newton}: $p_3=\sum x_k^3=\sum_{k\in C}\!\big[(\zeta^{3k}+\zeta^{-3k})+3x_k\big]$;
here $3k\bmod15$ runs over $\{3,6,12,6\}$ --- pairs of fifth
parts:
$\sum_{k\in C}(\zeta^{3k}+\zeta^{-3k})
=2\cdot2\cos\frac{2\pi}{5}+2\cdot2\cos\frac{4\pi}{5}
=(\sqrt5-1)-(1+\sqrt5)=-2$,
so $p_3=-2+3s_1=1$, and Newton's identity gives
$s_3=\frac{p_3-s_1p_2+s_2p_1}{3}=\frac{1-9-4}{3}=-4$
($p_2=\sum x_k^2=9$, $p_1=s_1=1$). Altogether
$P_{15}(x)=x^4-x^3-4x^2+4x+1$; it is the minimal polynomial
because $[K:\QQ]=4$ is already established.

\emph{The reflection.} The roots of $P_{30}$ are
$2\cos(\pi k/15)$ for $k\in\{1,7,11,13\}$ (the classes
$\{\pm k\}$ modulo $30$). For $m\in C$:
$-2\cos\frac{2\pi m}{15}=2\cos\!\big(\pi-\frac{2\pi m}{15}\big)
=2\cos\frac{\pi(15-2m)}{15}$, and $15-2m$ runs over
$\{13,11,7,1\}$ --- exactly the classes of $P_{30}$. Hence the
root set of $P_{30}$ is $\{-x\colon x\ \text{a root of}\ P_{15}\}$,
i.e. $P_{30}(x)=P_{15}(-x)=x^4+x^3-4x^2-4x+1$.

\emph{Corollary (the $b_{\mathrm{Ch}}$ forms).}
$b_{\mathrm{Ch}}(15)=1-\frac{x}{2}=\frac{7-\sqrt5-\beta}{8}$,
$b_{\mathrm{Ch}}(30)=1-\frac{x'}{2}=\frac{9-\sqrt5-\beta'}{8}$.
\end{proof}

\begin{theorem}[the exact layer: an integer certificate]\label{th:exact}
\emph{(a)} The element $x=m/4$ with $m=1+\sqrt5+\beta$ satisfies
$P_{15}$ identically: the expansion of $256\,P_{15}(m/4)$ in the
tower $\ZZ[\sqrt5][\beta]$ ($\beta^2=30-6\sqrt5$ --- an integer
element!) collapses to the zero quadruple $(0,0,0,0)$ --- a
\emph{pure integer identity} without a single floating-point
operation; likewise $256\,P_{30}((\sqrt5-1+\beta')/4)=0$ in
$\ZZ[\sqrt5][\beta']$. \emph{(b)} Both quartics agree modulo $2$:
$P_{15}\equiv P_{30}\equiv x^4+x^3+1\pmod2$ --- an irreducible
polynomial over $\mathbf{F}_2$ (no roots in $\mathbf{F}_2$;
division by the only irreducible quadratic $x^2+x+1$ leaves a
nonzero remainder) --- the machine fingerprint of the degree.
\emph{(c)} Isolation: the roots of $P_{15}$ are
$1.827091\ldots,\ 1.338261\ldots,\ -0.209057\ldots,\ -1.956295\ldots$
--- the largest root singles out the principal value
$2\cos(2\pi/15)$ uniquely.
\end{theorem}

\begin{proof}
(a) --- direct expansion: the tower arithmetic of
$\ZZ[\sqrt5][\beta]$ multiplies pairs $(A,B)$
($A+B\beta$ with $A,B\in\ZZ[\sqrt5]$) by the rule $\beta^2=D$
with an integral element $D=30\mp6\sqrt5$; all intermediate
results are integral because $\beta^2\in\ZZ[\sqrt5]$ and all
coefficients are integral. The laboratory certificate
(\texttt{\_bch\_exact\_residual}) returns the zero tower element
--- that is the identity. (b) --- reduction of coefficients:
$(1,-1,-4,4,1)\bmod2=(1,1,0,0,1)$ and
$(1,1,-4,-4,1)\bmod2=(1,1,0,0,1)$; the irreducibility of
$x^4+x^3+1$ over $\mathbf{F}_2$ is standard. (c) --- the numeric
check at dps 60: the radical forms match $2\cos$ with relative
error $0$ within precision; the ordering of the roots is strict.
\end{proof}

\begin{remark}
The quartics $P_{15}$, $P_{30}$ are norms of the cyclic quartic
field of discriminant $1125$; their reflected pair
$P_{30}(x)=P_{15}(-x)$ is the arithmetic echo of the fact that
the levels $15$ and $30$ share one cyclotomic field
$\QQ(\zeta_{15})=\QQ(\zeta_{30})$ (monograph 16/20, theorem H1).
The constructibility of the level (Gauss--Wantzel) and the
cyclicity $C_4$ are two faces of the same coin: the tower
$\QQ\subset\QQ(\sqrt5)\subset K$ has exactly one intermediate
floor.
\end{remark}

\begin{databox}[The exact layer v1.2 --- the numbers]
Radicals: $b_{\mathrm{Ch}}(15)=\frac{7-\sqrt5-\sqrt{30-6\sqrt5}}{8}
=0.086454542357399104497872428014682822\ldots$;
$b_{\mathrm{Ch}}(30)=\frac{9-\sqrt5-\sqrt{30+6\sqrt5}}{8}
=0.021852399266194362071433252130400468\ldots$
(the match with $1-\cos$ at dps 60 is $0$ within precision).
Minimal quartics: $x^4-x^3-4x^2+4x+1$ and
$x^4+x^3-4x^2-4x+1$; the reflection $P_{30}(x)=P_{15}(-x)$;
the tower residual $=0$ (an integer identity); $\mathrm{GF}(2)$:
both reduce to $x^4+x^3+1$ (irreducible); the pairing
$\beta\beta'=12\sqrt5$; $\beta^2+\beta'^2=60$; the field
$K=\QQ(\zeta_{15})^+$, $\mathrm{Gal}\cong C_4$,
$\disc K=1125=3^2\cdot5^3=5\cdot15\cdot15$ --- the same number
$1125$ as the volume $\mathrm{vol}_h$ of monograph 16/20.
\end{databox}

\begin{statusbox}[summary]
The exact layer of v1.2 is raised to the level of theorems: the
closed radicals of both constructible rungs, the minimal quartics
derived from first principles (Ramanujan sums + the tower norm +
Newton + the reflection), the Galois structure $C_4$ with the
unique quadratic floor $\QQ(\sqrt5)$, the integer tower
certificate and the $\mathrm{GF}(2)$ fingerprint. One field ---
two rungs.
\end{statusbox}

\begin{verifybox}[how to reproduce]
\texttt{python3 laboratory.py} $\to$ ``Designer $\to$ item 5''
(radicals + exact/numeric verdicts, $N=15/30$); batch mode:
a scenario with the \texttt{bch} items ($N=15$, $N=30$);
tests: \texttt{pytest tests/ -k bch}; Lean: the
\texttt{bch\_radicals} block; \texttt{--check-baseline}.
\end{verifybox}

\vfill
{\color{linkblue}\hrule height 0.6pt}
\vspace{0.5em}
{\small\color{gray}
License: individual exclusive license of Isaev Iskhak Khamzatovich.
All rights to the computations, formulas and texts belong to the author.
Verification: \texttt{python3 laboratory.py} (``Designer'', batch, tests).
}
\end{document}
